Since the observer is at the pole, the affect of the earth's rotation on the
transit could be neglected.

Then the Sun's angle at the earth extends as
\[
\theta_0 = \arcsin\!\left(\frac{2r_{\mathrm{sun}}}{1\,\mathrm{AU}}\right)
 \approx 32.0';
\]
the angular velocity of the Venus around the Sun, respected to the earth is
$\omega_1$,
\[
\omega_1 = \omega_{\mathrm{venus}} - \omega_{\mathrm{earth}}
 = \frac{2\pi}{T_{\mathrm{venus}}} - \frac{2\pi}{T_{\mathrm{earth}}}
 \approx 4.29\times 10^{-4}\ ('/\mathrm{s})
\]
\hfill(2 points)

For the observer on earth, Venus moved $\theta$ during the whole transit. Let
OE be perpendicular to AB,
\[
OA = 16', \quad \angle AOB = 90^\circ, \quad MN \parallel AB,
\]
So $OE = 11.3'$,
$OC = \dfrac{d'_{\mathrm{venus}}}{2} + \dfrac{r'_{\mathrm{sun}}}{2}$,
$d'_{\mathrm{venus}}$ is the angular size of Venus seen from Earth.
% FIGURE NEEDED: solar disk with transit chord geometry (points M, N, O, A,
% B, C, D, E, east/west directions) and auxiliary construction of Venus's
% orbit seen from Earth (O, C, D, F); 2010 proceedings, solutions of
% problem 14, p. 38.
\begin{align*}
d'_{\mathrm{venus}} &= \frac{2\times 0.949\times 6378}
                           {(1-0.723)\times 1\,\mathrm{AU}} \approx 1',\\
OC &\approx 16.5', \quad CD \approx 24.0',\\
CE &= \sqrt{OC^2 - OE^2} \approx 12.0'\\
CD &= 2CE = 24.0'
\end{align*}
So, $\theta = \angle CFD = 24.0'$, \hfill(3 points)

As shown on the picture, $\theta' = \angle COD$ is the additional angle that
Venus covered during the transit,
\[
\frac{\operatorname{tg}\dfrac{\theta}{2}}{\operatorname{tg}\dfrac{\theta'}{2}}
 = \frac{0.723}{(1-0.723)}, \quad
\operatorname{tg}\frac{\theta}{2} = \operatorname{tg} 12', \quad
\theta' = 9.195';
\]
\hfill(3 points)
\[
t_{\mathrm{transit}} = \frac{\theta'}{\omega_1}
 = \frac{9.195'}{4.29\times 10^{-4}\ '/\mathrm{s}}\times\cos\varepsilon,
 \text{ that is } 5^{\mathrm{h}}56^{\mathrm{m}}36^{\mathrm{s}},
\]
% sic: the factor cos(epsilon) appears in the source without epsilon being
% defined anywhere in the solution.
So the transit will finish at about $14^{\mathrm{h}}57^{\mathrm{m}}$.
\hfill(2 points)
